
\subsubsection{6.1 Why HumanEval-only Training Generalizes to Both Benchmarks}

Both HumanEval+ and MBPP+ evaluate Python functions that must pass test cases provided by the benchmark. The core skill required --- writing Python code that satisfies a functional specification --- is structurally similar across benchmarks despite surface-level differences in problem style. HumanEval problems are framed as algorithmic function completions with doctest-style specifications; MBPP problems present utility scripts with input/output pairs. Both require producing syntactically correct, functionally accurate Python.

HumanEval-only SFT may develop this general function-writing skill more effectively than MBPP-only SFT because HumanEval's doctest-driven training format more directly rewards producing a correct implementation rather than reproducing a particular coding style. MBPP-only SFT may develop a narrower ''utility script'' pattern that does not generalize as effectively.

An additional contributing factor: the MBPP-only training condition used only 120 problems (from the sanitized split, which loaded fewer samples than expected). This reduced dataset size may have limited MBPP-only's specialization potential relative to HumanEval-only's 164 problems. The direction of the finding --- HumanEval-only exceeding MBPP-only on MBPP+ --- is consistent across all three seeds with approximately 1.5--2.1pp margin, suggesting a real effect rather than a size-driven artifact. However, rerunning MBPP-only SFT with the full available MBPP training split would be necessary to rule out dataset size as a contributing factor.

\subsubsection{6.2 Equal-Mix Underperformance at 1.3B}

Equal-mix training produces 10.2\% HumanEval+ pass@1, substantially below both HumanEval-only (35.0\%) and MBPP-only (27.6\%), despite receiving the same total token budget. At 7B scale, equal-mix (39.6\%) matches or exceeds HumanEval-only (38.6\%), suggesting that this source specificity advantage of single-source training diminishes or reverses as model scale and pretraining coverage increase.

One interpretation: at 1.3B scale with limited SFT data, gradient updates from heterogeneous sources partially interfere, producing weaker specialization than focused single-source training. At 7B scale, the model's richer pretraining representations may provide a stronger foundation that is less sensitive to gradient direction during SFT, allowing equal-mix to produce broadly competent performance without the dilution penalty.

\subsubsection{6.3 The Role of Distributional Alignment}

The CodeBERT rank correlation result ($\rho$=1.0, p=0.042) shows that, among the four conditions studied, the condition with the highest embedding similarity to the test benchmark also achieves the highest pass@1, and this ordering is perfectly preserved. The mechanistic interpretation is that SFT gradient updates bias model behavior toward the distributional properties of the training corpus, and test-time performance benefits from alignment between training and test distributions.

This interpretation is qualified by three limitations. First, the statistical claim rests on n=4 conditions; the minimum achievable p-value equals the observed p-value, leaving no margin. Second, the alignment test is restricted to HumanEval+; MBPP+ data is insufficient for a parallel analysis. Third, the finding for HumanEval+ does not generalize straightforwardly to MBPP+: HumanEval-only also has the highest embedding similarity to MBPP+ (among the conditions tested), so the alignment predictor does not contradict the MBPP+ finding but also does not uniquely explain it.

\subsubsection{6.4 Scale Effects and the $\eta^2$ Limitation}

The near-deterministic behavior at 7B scale --- where each condition produces essentially identical results across three random seeds --- is itself an informative finding. At 1.3B, seed variance ($\sigma ^{2}\approx 0.017)$ is substantial: HumanEval-only produces 39.6\%, 25.6\%, 39.6\% across seeds 42, 123, 777, indicating sensitivity to initialization. At 7B, seed variance ($\sigma ^{2}\approx 0.0004)$ is negligible. This suggests that the relationship between source identity and benchmark performance is more stably encoded at larger scale, though whether this reflects reduced sensitivity to gradient noise, more robust pretraining representations, or both is not determined by these experiments.

The practical consequence for methodology: $\eta^2$ is not an appropriate metric for comparing source identity effect sizes across scales when seed sensitivity changes with scale. Researchers comparing fine-tuning effects across model scales should report absolute condition spread alongside or instead of $\eta^2$.

\subsubsection{6.5 Limitations}

\textbf{MBPP+ evaluation incomplete.} The full 4-condition MBPP+ transfer matrix is not available. H-E2's CSV output contains only HumanEval benchmark data (MBPP+ evaluation was not saved); H-C1 MBPP+ evaluation produced 4 of 12 parse failures. MBPP+ claims are restricted to the HumanEval-only vs MBPP-only comparison from H-M1. The existing 24 SFT checkpoints (1.3B and 7B) could be re-evaluated on MBPP+ without additional training.

\textbf{Statistical fragility of mechanistic claim.} With n=4 conditions, $\rho$=1.0 achieves the minimum permutation p-value (1/24 $\approx$ 0.042). A single rank swap eliminates significance. Dual-encoder concordance (MiniLM $\rho$=0.800 in the same direction) and mechanistic plausibility provide supporting evidence, but the mechanistic claim should not be characterized as robustly established.

\textbf{MBPP-only training set size.} The MBPP-only condition trained on 120 problems rather than the expected ~374. This may reduce MBPP-only's capacity for specialization, partially confounding source identity with dataset size.

\textbf{Format normalization as a potential confound.} A standardized template was applied across all conditions. MBPP-only SFT using the native MBPP problem format (which includes explicit input/output examples) might produce different MBPP+ performance; this was not tested.

\textbf{LeetCode training stability.} LeetCode-only SFT produced 0.0\% HumanEval+ pass@1 on 2/3 seeds. Training loss curves were not separately verified for LeetCode conditions. Training collapse cannot be ruled out as a contributing factor to the near-zero performance, though the distributional misalignment explanation (CodeBERT similarity 0.909, the lowest among sources) is mechanistically consistent.

\textbf{Scope.} All results are restricted to DeepSeek-Coder at 1.3B and 7B scales, Python code, and the specific benchmark pair (HumanEval+, MBPP+). Generalization to other languages, benchmarks, model families, or larger-scale SFT regimes is not established.

